\let\negmedspace\undefined
\let\negthickspace\undefined
\documentclass[journal]{IEEEtran}
\usepackage[a5paper, margin=10mm, onecolumn]{geometry}
%\usepackage{lmodern} % Ensure lmodern is loaded for pdflatex
\usepackage{tfrupee} % Include tfrupee package

\setlength{\headheight}{1cm} % Set the height of the header box
\setlength{\headsep}{0mm}     % Set the distance between the header box and the top of the text

\usepackage{gvv-book}
\usepackage{gvv}
\usepackage{cite}
\usepackage{amsmath,amssymb,amsfonts,amsthm}
\usepackage{algorithmic}
\usepackage{graphicx}
\usepackage{textcomp}
\usepackage{xcolor}
%\usepackage{txfonts}
\usepackage{listings}
\usepackage{enumitem}
\usepackage{mathtools}
\usepackage{gensymb}
\usepackage{comment}
\usepackage[breaklinks=true]{hyperref}
\usepackage{tkz-euclide} 
\usepackage{listings}
% \usepackage{gvv}                                        
\def\inputGnumericTable{}                                 
\usepackage[latin1]{inputenc}                                
\usepackage{color}                                            
\usepackage{array}                                            
\usepackage{longtable}                                       
\usepackage{calc}                                             
\usepackage{multirow}                                         
\usepackage{hhline}                                           
\usepackage{ifthen}                                           
\usepackage{lscape}
\usepackage{circuitikz}
\tikzstyle{block} = [rectangle, draw, fill=blue!20, 
    text width=4em, text centered, rounded corners, minimum height=3em]
\tikzstyle{sum} = [draw, fill=blue!10, circle, minimum size=1cm, node distance=1.5cm]
\tikzstyle{input} = [coordinate]
\tikzstyle{output} = [coordinate]


\begin{document}

\bibliographystyle{IEEEtran}
\vspace{3cm}

\title{5.13.16}
\author{AI25BTECH11030 -Sarvesh Tamgade}
{\let\newpage\relax\maketitle}

\renewcommand{\thefigure}{\theenumi}
\renewcommand{\thetable}{\theenumi}
\setlength{\intextsep}{10pt} 


\numberwithin{equation}{enumi}
\numberwithin{figure}{enumi}
\renewcommand{\thetable}{\theenumi}


\textbf{Question}:
The equation of a circle with origin as centre and passing through the vertices of an equilateral triangle whose median is of length $3a$ is
\begin{enumerate}
\item $x^2 + y^2 = 9a^2$
\item $x^2 + y^2 = 16a^2$
\item $x^2 + y^2 = 4a^2$
\item $x^2 + y^2 = a^2$
\end{enumerate}


\textbf{Solution}:  
\solution
Let the vertices of the equilateral triangle be $\vec{A}$, $\vec{B}$, and $\vec{C}$.And the radius of circle be r.

Since the circle is centered at the origin, and the radius is r:
\begin{align}
\vec{u} = -\myvec{0 \\ 0} = \myvec{0 \\ 0}
\end{align}
\begin{align}
f=-r^2
\end{align}

Since the circle passes through $\vec{A}$:
\begin{align}
\norm{\vec{x}^\top\vec{x}} + 2\vec{u}^T\vec{x} + f &= 0\\
\norm{\vec{x}^\top\vec{x}}+ 2\myvec{0 & 0}\vec{x} + f &= 0\\
\norm{\vec{x}^\top\vec{x}} - r^2 &= 0\\
\end{align}

For an equilateral triangle inscribed in a circle, let us place the vertices as:
\begin{align}
\vec{A} &= \myvec{r \\ 0}\\
\vec{B} &= \myvec{r\cos\theta \\ r\sin\theta}\\
\vec{C} &= \myvec{r\cos\theta \\ -r\sin\theta}
\end{align}

The median from vertex $\vec{A}$ goes to the midpoint of $\vec{B}$ and $\vec{C}$:
\begin{align}
\text{Midpoint} &= \frac{\vec{B} + \vec{C}}{2} = \myvec{r\cos\theta \\ 0}
\end{align}

The median vector is:
\begin{align}
\vec{m} &= \myvec{r\cos\theta \\ 0} - \myvec{r \\ 0} = \myvec{r(\cos\theta - 1) \\ 0}
\end{align}

Given that the median length is $3a$:
\begin{align}
\norm{\vec{m}^\top\vec{m}} &= (3a)^2 = 9a^2\\
r^2(\cos\theta - 1)^2 &= 9a^2
\end{align}

For an equilateral triangle, $\cos\theta = \cos\left(\frac{2\pi}{3}\right) = -\frac{1}{2}$:
\begin{align}
(\cos\theta - 1) &= -\frac{3}{2}\\
r^2  \frac{9}{4} &= 9a^2\\
r^2 &= 4a^2
\end{align}



Therefore, the equation of the circle is:
\begin{align}
\norm{\vec{x}^\top\vec{x}} + 2\myvec{0 & 0}\vec{x} - 4a^2 &= 0\\
\norm{\vec{x}^\top\vec{x}} - 4a^2 &= 0
\end{align}

\textbf{Answer:} The orthocentre of the triangle is
Therefore, the equation of the circle is:
\begin{align}
x^2 + y^2 = 4a^2
\end{align}

The answer is (c).

\newpage
\begin{figure}[htbp]
    \centering
    \includegraphics[width=0.65\linewidth]{FIG/fig1.png}
    \caption{Vector Representation}
    \label{fig:FIG/fig1.png}
    \end{figure}

\end{document}  